\subsection{Steerability}
\label{sec:steerable}
\label{sec:steering}
We now shift our focus from measuring the default alignment of LM opinions with those of various demographics groups \emph{without} prompting, to studying their steerability \emph{with} group-specific prompting. This is especially important in settings such as personalization, where
a key measure of performance is an LM's ability to adapt to represent the opinion of various demographic groups.

\paragraph{The metric.} We measure steerability as the average opinion alignment, across dataset questions, between an LM and a particular demographic group $G$ -- where the model is prompted with group information in its context.
Since our goal is to test whether a model \emph{can} be steered toward a group, we consider three prompting strategies---\textsc{QA},\textsc{BIO},\textsc{PORTRAY} (see Section~\ref{sec:prompting})---for each question and choose the one that works best.
Concretely, we measure steerability as:

\begin{equation}
    \steer^{G}_m(Q) =  \frac{1}{|Q|} \sum_{q \in Q} \max_{c_G \in \small{[\textsc{QA},\textsc{BIO},\textsc{POR}]}} \simi(D_m(q; c_G), D_{G}(q))  \nonumber
\end{equation}
where $D_m(q; c_G)$ denotes the LM opinion distribution conditioned on the group-specific context $c_G$.
A higher $\steer^{G}_m$ score indicates that the model is better aligned to the opinions of the given group.
Note that unlike default subgroup representativeness, an LM's steerability could be simultaneously high for multiple (disagreeing) groups.
In fact, in many cases, we might want disparities in the default subgroup representativeness scores of an LM to be remedied by steering.


\paragraph{Steering does not solve opinion misalignment.}

We attempt to steer current LMs towards one of 22 demographic groups (\eg, Republican, Asian, Jewish) listed in Appendix Table~\ref{tab:steer_groups} on a subset $Q_S$ of 500 highly contentious questions from \bname.
In Figure~\ref{fig:steerability_overall}, we compare different LMs in terms of their ability to match the opinions of a subgroup on these questions, by default and with steering ($\steer^{G}_m(Q_S)$ from Section~\ref{sec:rep}).

Most LMs (with the exception of \model{ada}) do become somewhat more representative of a subpopulation post-steering.
However, none of the disparities in group opinion alignment of an LM disappear after steering, with \model{text-davinci-002} showing the smallest post-steering alignment gap across groups. In most cases, we see the representativeness of all groups improving by a constant factor---indicating that the LM still does better on some groups than others.
In Appendix Figure~\ref{figapp:steerability_group}, we visualize which LMs are most effective at adapting towards a particular group: \eg, \model{j1-grande-v2-beta} for Southerners and \model{text-davinci-002} for liberals.

